\section{Revised AI Interactions}

Seven additional ChatGPT interactions (October 19,
2026) were used to review and document the greedy
project. Each entry gives the exact prompt and a
summary of the response.

\begin{enumerate}

\item \textbf{Problem selection and abstraction.}
\emph{Prompt:} ``I am completing a graduate-level
Analysis of Algorithms project that requires a
real-world application of a greedy algorithm. I want
to investigate traffic signal scheduling to help avoid
overlapping green-light opportunities and potentially
reduce scheduling conflicts. Can you explain how to
formulate a simplified traffic signal scheduling
problem using interval scheduling? Define the candidate
green-light windows, starting and finishing times,
compatibility constraints, objective function,
assumptions, and limitations. Clearly distinguish
maximizing the number of compatible windows from
minimizing actual traffic congestion.''
\emph{Result:} Defined windows $w_i=(s_i,f_i)$, the
compatibility condition $f_i \le s_j$, and the
objective $\max|S|$; listed the model's assumptions;
and explained that it maximizes compatible windows
rather than minimizing congestion.

\item \textbf{Greedy algorithm and C++ implementation.}
\emph{Prompt:} ``For my COT5405 Analysis of Algorithms
project, I modeled traffic signal scheduling as a
maximum-cardinality interval scheduling problem. Each
candidate green-light window has an ID, start time, and
finish time. Explain the earliest-finish-time greedy
algorithm step by step, provide pseudocode, and show a
complete C++17 implementation using a struct, vector,
and std::sort. Use the example windows (1,0,10),
(2,2,5), (3,5,8), (4,7,12), and (5,8,14). Explain why
the selected windows are 2, 3, and 5.''
\emph{Result:} Provided pseudocode, a C++17
implementation, a walkthrough selecting windows 2, 3,
and 5, and an $O(n\log n)$ analysis.

\item \textbf{Code review and debugging.}
\emph{Prompt:} ``I have written a C++17 greedy algorithm
to select the maximum number of mutually compatible
green-light windows. Review my implementation for
correctness, compilation errors, sorting logic,
boundary conditions, and output formatting. Explain
every necessary correction and why it matters. Preserve
my original program structure as much as possible
instead of rewriting the entire implementation. Also
identify useful test cases, including overlapping
intervals, back-to-back intervals, empty input, and
intervals with identical finishing times. Here is my
implementation:'' [followed by the complete 25-line
\texttt{greedy.cpp} program, omitted here for space]
\emph{Result:} Confirmed the code was correct for the
example; noted that initializing \texttt{lastFinish} to
$-1$ assumes nonnegative start times and suggested a
generalization; proposed edge-case tests.

\item \textbf{Experimental runtime analysis.}
\emph{Prompt:} ``I need to experimentally verify the
O(n log n) running-time complexity of my C++
earliest-finish-time greedy interval scheduling
algorithm. Help me design and implement a C++17
benchmark that generates random green-light windows for
input sizes 1,000, 5,000, 10,000, 25,000, 50,000,
100,000, and 200,000. Use a fixed random seed for
reproducibility, run 10 repetitions per input size,
measure sorting and greedy selection using
std::chrono::steady\_clock, exclude input generation
from the timed section, compute median execution time
in milliseconds, and export InputSize, MedianTimeMs, and
NLogN to Greedy/results.csv. Explain compilation using
g++ -std=c++17 -O2 and the limitations of runtime
measurements.''
\emph{Result:} Provided the benchmark in
Listing~\ref{lst:experiment} (seed 42, median of ten
runs, CSV export) and discussed measurement noise.

\item \textbf{Runtime graph generation.}
\emph{Prompt:} ``I have a CSV file at Greedy/results.csv
containing the columns InputSize, MedianTimeMs, and
NLogN from my C++ greedy algorithm benchmark. Show me
how to write and execute Greedy/plot\_results.py using
Python, csv.DictReader, and Matplotlib. Plot measured
median execution time against input size and compare it
with a scaled theoretical n log2(n) curve. Save the
figure as Greedy/runtime\_graph.png. Explain the meaning
of the graph, how it supports the theoretical runtime
analysis, and how to diagnose incorrect CSV headers or
missing files.''
\emph{Result:} Provided the plotting script
(Listing~\ref{lst:plot}), explained scaling the
$n\log_2 n$ curve to the largest measurement, and
described common file and header errors.

\item \textbf{Correctness proof.}
\emph{Prompt:} ``I need a formal proof of correctness for
the earliest-finish-time greedy algorithm used in my
traffic signal scheduling project. The objective is to
maximize the number of nonoverlapping candidate
green-light windows, where each window has a start and
finish time and two consecutive windows are compatible
if the next starts at or after the previous finishes.
Provide a clear theorem and an exchange-argument proof
showing why replacing the first interval of an optimal
schedule with the earliest-finishing interval preserves
optimality. Explain why the argument can be applied
repeatedly to the remaining intervals. Provide the
proof in LaTeX and state its assumptions.''
\emph{Result:} Gave a theorem and exchange-argument
proof, completed by induction on the number of windows.

\item \textbf{Report preparation and AI disclosure.}
\emph{Prompt:} ``I am preparing a COT5405 algorithms
group project report in Overleaf using the IEEE
conference template. My portion covers greedy traffic
signal window scheduling. Help me organize the report
into a real-world problem description, mathematical
abstraction, algorithm and pseudocode, correctness
proof, theoretical O(n log n) analysis, empirical
runtime table and graph, bibliography, and source-code
appendix. Explain how to divide sections into separate
LaTeX files using input commands, include a PNG figure,
include C++ and Python code using the listings package,
and use BibTeX. Also help me prepare a transparent AI-use
appendix distinguishing exact original prompts and
responses from summaries or newly generated
interactions. My partner's divide-and-conquer section
will be added later.''
\emph{Result:} Proposed the IEEE report structure and
\texttt{\textbackslash input} file layout, showed how to
include figures, code listings, and BibTeX references,
and advised keeping original and revised AI records
separate.

\end{enumerate}
